% Complete reviewed Unit 8; editorial master: chapters/ch08/answer-key.md
% Markdown SHA-256: 41f74e8c6e77abdc6b0f491d4efb95a214d49d4ed60a3572280d9802d562e4f6
% One-use initial assembly followed by complete content and page review.
\newcommand{\runningtitle}{Unit 8. Answer key}
% Layout inherited from the reviewed previous course; set a running title per unit.
% This is a small hand-maintained TeX layout, not a conversion pipeline.
\RequirePackage{fix-cm}
\documentclass[11pt,twoside]{article}
\usepackage{fontspec}
\setmainfont{Cambria}[
  Renderer=HarfBuzz,
  BoldFont=cambriab.ttf,
  ItalicFont=cambriai.ttf,
  BoldItalicFont=cambriaz.ttf
]
\setmonofont{lmmono10-regular.otf}[BoldFont=lmmonolt10-bold.otf]
\usepackage[paperwidth=165mm,paperheight=241mm,inner=21mm,outer=19mm,
  top=22mm,bottom=20mm,headheight=14pt,headsep=7mm,footskip=10mm]{geometry}
\usepackage{xcolor,array,booktabs,longtable,tabularx}
\usepackage{enumitem,fancyhdr}
\usepackage[unicode,hidelinks]{hyperref}
\definecolor{ink}{HTML}{000000}
\definecolor{muted}{HTML}{4C5558}
\setlength{\parindent}{0pt}
\setlength{\parskip}{3.5pt}
\AtBeginDocument{\fontsize{11}{14.3}\selectfont}
\setlength{\tabcolsep}{4pt}
\renewcommand{\arraystretch}{1.10}
\setlist[itemize]{leftmargin=1.3em,label=\textbullet,itemsep=2pt,topsep=3pt,parsep=0pt}
\setlist[enumerate]{leftmargin=1.8em,itemsep=3pt,topsep=4pt,parsep=0pt}
\widowpenalty=10000
\clubpenalty=10000
\emergencystretch=1em
\raggedbottom
\pagestyle{fancy}
\fancyhf{}
\providecommand{\runningtitle}{Old English}
\fancyhead[LE,RO]{\small\itshape \runningtitle}
\fancyhead[LO,RE]{\small Old English}
\fancyfoot[LE,RO]{\small\thepage}
\renewcommand{\headrulewidth}{0.3pt}
\newcommand{\booktitle}[2]{
  {\fontsize{9}{11}\selectfont\color{muted}#1\par}
  \vspace{3pt}
  {\fontsize{20}{24}\selectfont\color{ink}#2\par}
  \vspace{6pt}
}
\newcommand{\source}[1]{\par{\footnotesize\color{muted}#1\par}}
\newcommand{\readingtext}[1]{
  \par\addvspace{5pt}
  \begingroup\leftskip=1.5em
    \fontsize{12}{16}\selectfont #1\par
  \endgroup\addvspace{5pt}
}
\makeatletter
\renewcommand{\section}{\@startsection{section}{1}{\z@}
  {-2.4ex \@plus -0.6ex \@minus -0.2ex}{0.8ex \@plus 0.2ex}
  {\normalfont\fontsize{12}{15}\selectfont\bfseries\color{ink}}}
\renewcommand{\subsection}{\@startsection{subsection}{2}{\z@}
  {-1.8ex \@plus -0.5ex \@minus -0.2ex}{0.5ex \@plus 0.2ex}
  {\normalfont\fontsize{11}{14}\selectfont\bfseries\color{ink}}}
\makeatother

\hyphenpenalty=10000
\exhyphenpenalty=10000
\emergencystretch=3em
\newcommand{\itempara}[2]{\par\addvspace{3pt}\noindent\hangindent=1.8em\hangafter=1\makebox[1.8em][l]{#1.}#2\par}
\newcommand{\reserve}[1]{\par\penalty-100\begingroup\dimen0=\pagegoal\advance\dimen0 by-\pagetotal\ifdim\dimen0<#1\newpage\fi\endgroup}
\begin{document}
\addfontfeatures{Ligatures=TeXOff}
\raggedright
\booktitle{UNIT 8}{Answer key}

% mdparagraph:P001:begin
The models use the selected forms of this course. Other attested forms and word orders may be valid Old English. Assess their meaning and grammar separately from compliance with an expressly requested form or construction. Preserve the participants, number, tense and force of the task; do not judge an answer only by its resemblance to the model string.
% mdparagraph:P001:end

% mdparagraph:P002:begin
In paradigm tasks, give the selected table forms. \textbf{Hwām} is a documented dative alternative to selected \textbf{hwǣm}, not an accusative form. The ordinary relative pattern assigns case from the relative's role within its own clause. The lesson's recognition note on case attraction qualifies that rule; it does not replace the stated target in these compositions. The source reading's local helps remain available, as does \textbf{gesęte} in Cook I.8.13. Supplied forms and constructions are not tests of unaided recall. [1]
% mdparagraph:P002:end

\reserve{7\baselineskip}
% mdparagraph:P003:begin
\section*{A. Complete forms}
% mdparagraph:P003:end

\par\begin{minipage}{\linewidth}\raggedright
% item: NEW-U08-A01
% key: KEY-NEW-U08-A01
% mdparagraph:P004:begin
\textbf{A01. Who? What?}
% mdparagraph:P004:end

% mdtable:T01:begin
\par\begingroup
\fontsize{10.5}{13.2}\selectfont
\setlength{\LTleft}{0pt}\setlength{\LTright}{\fill}
\noindent
\begin{tabular}{@{}>{\raggedright\arraybackslash}p{35.73mm}>{\raggedright\arraybackslash}p{57.22mm}>{\raggedright\arraybackslash}p{25.43mm}@{}}
\toprule
\textbf{Case} & \textbf{Personal: masculine/\allowbreak{}feminine} & \textbf{Neuter} \\
\midrule
Nominative singular & hwā & hwæt \\
Genitive singular & hwæs & hwæs \\
Dative singular & hwǣm & hwǣm \\
Accusative singular & hwone & hwæt \\
\bottomrule
\end{tabular}
\par\endgroup

% mdtable:T01:end
% mdparagraph:P005:begin
\textbf{Target:} the complete singular nominative, genitive, dative and accusative system. The personal series serves both masculine and feminine referents; there is no separate feminine or plural paradigm. \textbf{Hwām} is a valid dative variant. Instrumental \textbf{hwȳ/\allowbreak{}hwon} belongs to the separate reference display and does not replace personal accusative \textbf{hwone}. Nominative \textbf{hwā} cannot supply every English “who/\allowbreak{}whom” function.
% mdparagraph:P005:end

\end{minipage}\par\addvspace{4pt}

\par\begin{minipage}{\linewidth}\raggedright
% item: NEW-U08-A02
% key: KEY-NEW-U08-A02
% mdparagraph:P006:begin
\textbf{A02. Hwilc: which?}
% mdparagraph:P006:end

% mdtable:T02:begin
\par\begingroup
\fontsize{10.5}{13.2}\selectfont
\setlength{\LTleft}{0pt}\setlength{\LTright}{\fill}
\noindent
\begin{tabular}{@{}>{\raggedright\arraybackslash}p{36.86mm}>{\raggedright\arraybackslash}p{26.23mm}>{\raggedright\arraybackslash}p{26.23mm}>{\raggedright\arraybackslash}p{26.23mm}@{}}
\toprule
\textbf{Requested cell} & \textbf{Masculine} & \textbf{Neuter} & \textbf{Feminine} \\
\midrule
Genitive singular & hwilces & hwilces & hwilcre \\
Dative singular & hwilcum & hwilcum & hwilcre \\
Accusative singular & hwilcne & hwilc & hwilce \\
Nominative plural & hwilce & hwilc & hwilca \\
\bottomrule
\end{tabular}
\par\endgroup

% mdtable:T02:end
% mdparagraph:P007:begin
\textbf{Target:} strong agreement in case, number and gender. Feminine genitive and dative singular coincide, but masculine accusative singular \textbf{hwilcne} differs from the zero-ending neuter. Neuter nominative plural is \textbf{hwilc}, not an assumed \textbf{-e} form copied from the masculine.
% mdparagraph:P007:end

\end{minipage}\par\addvspace{4pt}

\par\begin{minipage}{\linewidth}\raggedright
% item: NEW-U08-A03
% key: KEY-NEW-U08-A03
% mdparagraph:P008:begin
\textbf{A03. Which of two? Of what sort?}
% mdparagraph:P008:end

% mdtable:T03:begin
\par\begingroup
\fontsize{10.5}{13.2}\selectfont
\setlength{\LTleft}{0pt}\setlength{\LTright}{\fill}
\noindent
\begin{tabular}{@{}>{\raggedright\arraybackslash}p{56.96mm}>{\raggedright\arraybackslash}p{30.71mm}>{\raggedright\arraybackslash}p{30.71mm}@{}}
\toprule
\textbf{Requested cell} & \textbf{hwæðer} & \textbf{hūlic} \\
\midrule
Masculine accusative singular & hwæðerne & hūlicne \\
Feminine dative singular & hwæðerre & hūlicre \\
Neuter nominative plural & hwæðer & hūlic \\
Genitive plural, all genders & hwæðerra & hūlicra \\
\bottomrule
\end{tabular}
\par\endgroup

% mdtable:T03:end
% mdparagraph:P009:begin
\textbf{Target:} each word's selected strong forms. \textbf{Hwæðer} has contracted forms elsewhere, such as \textbf{hwæðres/\allowbreak{}hwæðrum}; that does not license deleting its second vowel in every cell. In particular, distinguish \textbf{hwæðerne} from \textbf{hwæðerre}. Selected neuter plural \textbf{hwæðer} has no \textbf{-u}; do not copy the \textbf{-u} of \textbf{ōðru} mechanically. Documented variant spellings remain distinct from the requested selected table. Cook's \textbf{hūlic} keeps the length mark on \textbf{ū}.
% mdparagraph:P009:end

\end{minipage}\par\addvspace{4pt}

\par\begin{minipage}{\linewidth}\raggedright
% item: NEW-U08-A04
% key: KEY-NEW-U08-A04
% mdparagraph:P010:begin
\textbf{A04. Ilca: the same.}
% mdparagraph:P010:end

% mdtable:T04:begin
\par\begingroup
\fontsize{10.5}{13.2}\selectfont
\setlength{\LTleft}{0pt}\setlength{\LTright}{\fill}
\noindent
\begin{tabular}{@{}>{\raggedright\arraybackslash}p{35.82mm}>{\raggedright\arraybackslash}p{46.74mm}>{\raggedright\arraybackslash}p{35.82mm}@{}}
\toprule
\textbf{Case} & \textbf{Masculine singular} & \textbf{Neuter plural} \\
\midrule
Nominative & ilca & ilcan \\
Genitive & ilcan & ilcra \\
Dative & ilcan & ilcum \\
Accusative & ilcan & ilcan \\
\bottomrule
\end{tabular}
\par\endgroup

% mdtable:T04:end
% mdparagraph:P011:begin
\textbf{Target:} the two complete requested portions of the weak declension. \textbf{Ilca} follows the weak system; the unit assigns no strong alternative. Masculine nominative singular \textbf{ilca} is not a general invariant form for “same.” The plural genitive and dative are \textbf{ilcra/\allowbreak{}ilcum}, not singular \textbf{ilcan} repeated in every row.
% mdparagraph:P011:end

\end{minipage}\par\addvspace{4pt}

\par\begin{minipage}{\linewidth}\raggedright
% item: NEW-U08-A05
% key: KEY-NEW-U08-A05
% mdparagraph:P012:begin
\textbf{A05. Self.}
% mdparagraph:P012:end

% mdtable:T05:begin
\par\begingroup
\fontsize{10.5}{13.2}\selectfont
\setlength{\LTleft}{0pt}\setlength{\LTright}{\fill}
\noindent
\begin{tabular}{@{}>{\raggedright\arraybackslash}p{36.86mm}>{\raggedright\arraybackslash}p{26.23mm}>{\raggedright\arraybackslash}p{26.23mm}>{\raggedright\arraybackslash}p{26.23mm}@{}}
\toprule
\textbf{Nominative singular} & \textbf{Masculine} & \textbf{Neuter} & \textbf{Feminine} \\
\midrule
Strong & self & self & self \\
Weak & selfa & selfe & selfe \\
\bottomrule
\end{tabular}
\par\endgroup

% mdtable:T05:end
% mdparagraph:P013:begin
The strong masculine accusative singular is \textbf{selfne}. The strong dative plural, in all three genders, is \textbf{selfum}.
% mdparagraph:P013:end

% mdparagraph:P014:begin
\textbf{Target:} distinguish strong from weak agreement and nominative from oblique case. Identical strong nominative forms do not make \textbf{self} indeclinable. A masculine accusative referent requires \textbf{selfne} in the selected strong pattern, not nominative \textbf{self}.
% mdparagraph:P014:end

\end{minipage}\par\addvspace{4pt}

\par\begin{minipage}{\linewidth}\raggedright
% item: NEW-U08-A06
% key: KEY-NEW-U08-A06
% mdparagraph:P015:begin
\textbf{A06. Some, any, each, such.}
% mdparagraph:P015:end

% mdtable:T06:begin
\par\begingroup
\fontsize{10.5}{13.2}\selectfont
\setlength{\LTleft}{0pt}\setlength{\LTright}{\fill}
\noindent
\begin{tabular}{@{}>{\raggedright\arraybackslash}p{56.96mm}>{\raggedright\arraybackslash}p{30.71mm}>{\raggedright\arraybackslash}p{30.71mm}@{}}
\toprule
\textbf{Requested cell} & \textbf{sum} & \textbf{ǣnig} \\
\midrule
Masculine accusative singular & sumne & ǣnigne \\
Feminine genitive singular & sumre & ǣnigre \\
Neuter nominative plural & sumu & ǣnigu \\
Dative plural, all genders & sumum & ǣnigum \\
\bottomrule
\end{tabular}
\par\endgroup

% mdtable:T06:end
% mdtable:T07:begin
\par\begingroup
\fontsize{10.5}{13.2}\selectfont
\setlength{\LTleft}{0pt}\setlength{\LTright}{\fill}
\noindent
\begin{tabular}{@{}>{\raggedright\arraybackslash}p{56.96mm}>{\raggedright\arraybackslash}p{30.71mm}>{\raggedright\arraybackslash}p{30.71mm}@{}}
\toprule
\textbf{Requested cell} & \textbf{gehwilc} & \textbf{swilc} \\
\midrule
Masculine accusative singular & gehwilcne & swilcne \\
Feminine genitive singular & gehwilcre & swilcre \\
Neuter nominative plural & gehwilc & swilc \\
Dative plural, all genders & gehwilcum & swilcum \\
\bottomrule
\end{tabular}
\par\endgroup

% mdtable:T07:end
% mdparagraph:P016:begin
\textbf{Target:} the same case and agreement distinctions across four lexical stems. Selected \textbf{sumu/\allowbreak{}ǣnigu} contrast with zero-ending \textbf{gehwilc/\allowbreak{}swilc}. Analogical \textbf{ǣnig} without \textbf{-u} is a documented alternative, but \textbf{ǣnigu} is the selected paradigm answer. Retain \textbf{i} in the selected \textbf{ǣnig-} forms; do not invent contractions from another adjective's pattern. Displaying plural \textbf{gehwilc} does not by itself assign every distributive use of that plural.
% mdparagraph:P016:end

\end{minipage}\par\addvspace{4pt}

\par\begin{minipage}{\linewidth}\raggedright
% item: NEW-U08-A07
% key: KEY-NEW-U08-A07
% mdparagraph:P017:begin
\textbf{A07. Līf: life; an individual lived life.}
% mdparagraph:P017:end

% mdtable:T08:begin
\par\begingroup
\fontsize{10.5}{13.2}\selectfont
\setlength{\LTleft}{0pt}\setlength{\LTright}{\fill}
\noindent
\begin{tabular}{@{}>{\raggedright\arraybackslash}p{39.46mm}>{\raggedright\arraybackslash}p{39.46mm}>{\raggedright\arraybackslash}p{39.46mm}@{}}
\toprule
\textbf{Case} & \textbf{Singular} & \textbf{Plural} \\
\midrule
Nominative & līf & līf \\
Genitive & līfes & līfa \\
Dative & līfe & līfum \\
Accusative & līf & līf \\
\bottomrule
\end{tabular}
\par\endgroup

% mdtable:T08:end
% mdparagraph:P018:begin
\textbf{Target:} a complete neuter singular and plural declension, with the individual-life sense used in \textbf{ūru līf}, “our lives.” The plural refers to the several lives lived by the people denoted by \textbf{wē}. It is an ordinary contextual plural, not an unexplained exceptional plural of an invariably abstract noun. The zero ending leaves nominative and accusative plural identical to the corresponding singular forms; agreement and context decide the number.
% mdparagraph:P018:end

\end{minipage}\par\addvspace{4pt}

\par\begin{minipage}{\linewidth}\raggedright
% item: NEW-U08-A08
% key: KEY-NEW-U08-A08
% mdparagraph:P019:begin
\textbf{A08. Ætēowian: appear.}
% mdparagraph:P019:end

% mdtable:T09:begin
\par\begingroup
\fontsize{10.5}{13.2}\selectfont
\setlength{\LTleft}{0pt}\setlength{\LTright}{\fill}
\noindent
\begin{tabular}{@{}>{\raggedright\arraybackslash}p{56.24mm}>{\raggedright\arraybackslash}p{31.07mm}>{\raggedright\arraybackslash}p{31.07mm}@{}}
\toprule
\textbf{Indicative person and number} & \textbf{Present} & \textbf{Preterite} \\
\midrule
First singular & ætēowie & ætēowode \\
Second singular & ætēowast & ætēowodest \\
Third singular & ætēowað & ætēowode \\
First plural & ætēowiað & ætēowodon \\
Second plural & ætēowiað & ætēowodon \\
Third plural & ætēowiað & ætēowodon \\
\bottomrule
\end{tabular}
\par\endgroup

% mdtable:T09:end
% mdtable:T10:begin
\par\begingroup
\fontsize{10.5}{13.2}\selectfont
\setlength{\LTleft}{0pt}\setlength{\LTright}{\fill}
\noindent
\begin{tabular}{@{}>{\raggedright\arraybackslash}p{54.9mm}>{\raggedright\arraybackslash}p{31.74mm}>{\raggedright\arraybackslash}p{31.74mm}@{}}
\toprule
\textbf{Subjunctive number} & \textbf{Present} & \textbf{Preterite} \\
\midrule
Singular: persons 1, 2, 3 & ætēowie & ætēowode \\
Plural: persons 1, 2, 3 & ætēowien & ætēowoden \\
\bottomrule
\end{tabular}
\par\endgroup

% mdtable:T10:end
% mdparagraph:P020:begin
\textbf{Target:} the full assigned weak-II indicative and subjunctive systems. Each subjunctive row includes all three persons of its number; a dual subject also takes a plural finite verb. Present \textbf{-ie/\allowbreak{}-ien} contrasts with the preterite stem \textbf{ætēowod-}. Some forms coincide across moods: \textbf{ætēowie} is also first-singular present indicative, and \textbf{ætēowode} also first- or third-singular preterite indicative. Analyse the clause, not merely the spelling. This intransitive “appear” is distinct from the transitive verb “reveal, display.”
% mdparagraph:P020:end

\end{minipage}\par\addvspace{4pt}

\par\begin{minipage}{\linewidth}\raggedright
% item: NEW-U08-A09
% key: KEY-NEW-U08-A09
% mdparagraph:P021:begin
\textbf{A09. Shining; commands to bless.}
% mdparagraph:P021:end

% mdtable:T11:begin
\par\begingroup
\fontsize{10.5}{13.2}\selectfont
\setlength{\LTleft}{0pt}\setlength{\LTright}{\fill}
\noindent
\begin{tabular}{@{}>{\raggedright\arraybackslash}p{55.66mm}>{\raggedright\arraybackslash}p{31.36mm}>{\raggedright\arraybackslash}p{31.36mm}@{}}
\toprule
\textbf{Scīnan: subjunctive number} & \textbf{Present} & \textbf{Preterite} \\
\midrule
Singular: persons 1, 2, 3 & scīne & scine \\
Plural: persons 1, 2, 3 & scīnen & scinen \\
\bottomrule
\end{tabular}
\par\endgroup

% mdtable:T11:end
% mdtable:T12:begin
\par\begingroup
\fontsize{10.5}{13.2}\selectfont
\setlength{\LTleft}{0pt}\setlength{\LTright}{\fill}
\noindent
\begin{tabular}{@{}>{\raggedright\arraybackslash}p{83.32mm}>{\raggedright\arraybackslash}p{37.87mm}@{}}
\toprule
\textbf{Geblētsian: second-person command} & \textbf{Form} \\
\midrule
One person addressed & geblētsa \\
Two or more people addressed & geblētsiað \\
\bottomrule
\end{tabular}
\par\endgroup

% mdtable:T12:end
% mdparagraph:P022:begin
\textbf{Target:} retain long \textbf{ī} in the present subjunctive and short \textbf{i} in the preterite subjunctive, which uses the strong preterite plural stem. Do not substitute singular indicative \textbf{scān}. Each subjunctive entry covers persons 1, 2 and 3 in its number. The commands are second person; \textbf{geblētsiað} also coincides in spelling with a present indicative plural, but command force determines its use here.
% mdparagraph:P022:end

\end{minipage}\par\addvspace{4pt}

\reserve{7\baselineskip}
% mdparagraph:P023:begin
\section*{Q. Questions}
% mdparagraph:P023:end

\par\begin{minipage}{\linewidth}\raggedright
% item: NEW-U08-Q01
% key: KEY-NEW-U08-Q01
% mdparagraph:P024:begin
\textbf{Q01. Yes/\allowbreak{}no questions.}
% mdparagraph:P024:end

% mdparagraph:P025:begin
\itempara{1}{\textbf{Hæfð hē scip?} — Does he have a ship?}
% mdparagraph:P025:end

% mdparagraph:P026:begin
\itempara{2}{\textbf{Sindon ðā męnn gōde?} — Are the men good?}
% mdparagraph:P026:end

% mdparagraph:P027:begin
\itempara{3}{\textbf{Gesāwe ðū ðæt scip?} — Did you see that ship?}
% mdparagraph:P027:end

% mdparagraph:P028:begin
\textbf{Target:} finite verb before its subject, with third-singular present \textbf{hæfð}, third-plural present \textbf{sindon}, and second-singular preterite \textbf{gesāwe} retained. The question does not change these to infinitives or present tense. Another attested question order may be valid, but these models practise the ordinary finite-first pattern; no added \textbf{dōn} is required.
% mdparagraph:P028:end

\end{minipage}\par\addvspace{4pt}

\newpage
\par\begin{minipage}{\linewidth}\raggedright
% item: NEW-U08-Q02
% key: KEY-NEW-U08-Q02
% mdparagraph:P029:begin
\textbf{Q02. Who, whom, what.}
% mdparagraph:P029:end

% mdparagraph:P030:begin
\itempara{1}{\textbf{Hwā hæfð ðæt scip?} — \textbf{Hwā} is nominative, the subject of \textbf{hæfð}.}
% mdparagraph:P030:end

% mdparagraph:P031:begin
\itempara{2}{\textbf{Hwone siehð se mann?} — \textbf{Hwone} is accusative, the person seen; \textbf{se mann} is the nominative subject.}
% mdparagraph:P031:end

% mdparagraph:P032:begin
\itempara{3}{\textbf{Hwæt hæfð se mann?} — \textbf{Hwæt} is accusative neuter, the thing possessed; \textbf{se mann} is the subject.}
% mdparagraph:P032:end

% mdparagraph:P033:begin
\textbf{Target:} ask respectively who possesses the ship, whom the man sees, and what the man possesses. English question auxiliary “does” has no separate word in these models. \textbf{Hwæt} also has a nominative use, but its role here is object. \textbf{Hwǣm/\allowbreak{}hwām} would be dative, not the requested direct object “whom.”
% mdparagraph:P033:end

\end{minipage}\par\addvspace{4pt}

\par\begin{minipage}{\linewidth}\raggedright
% item: NEW-U08-Q03
% key: KEY-NEW-U08-Q03
% mdparagraph:P034:begin
\textbf{Q03. Recipient and possessor.}
% mdparagraph:P034:end

% mdparagraph:P035:begin
\itempara{1}{\textbf{Hwǣm sęlst ðū ðone fisc?} — To whom do you give that fish?}
% mdparagraph:P035:end

% mdparagraph:P036:begin
\itempara{2}{\textbf{Hwæs scip siehst ðū?} — Whose ship do you see?}
% mdparagraph:P036:end

% mdparagraph:P037:begin
\textbf{Target:} \textbf{hwǣm} is dative as the recipient of \textbf{sęllan}; \textbf{ðone fisc} is the accusative thing given. \textbf{Hwæs} is genitive as possessor of \textbf{scip}, while the whole phrase \textbf{hwæs scip} is the accusative object of \textbf{siehst}. Both finite verbs are second-singular present. \textbf{Hwām} is a valid dative alternative. Do not make the possessor accusative merely because its noun phrase is an object.
% mdparagraph:P037:end

\end{minipage}\par\addvspace{4pt}

\par\begin{minipage}{\linewidth}\raggedright
% item: NEW-U08-Q04
% key: KEY-NEW-U08-Q04
% mdparagraph:P038:begin
\textbf{Q04. Which of two? What sort?}
% mdparagraph:P038:end

% mdparagraph:P039:begin
\itempara{1}{\textbf{Hwæðer mann hæfð ðæt scip?} — Which of the two men has that ship?}
% mdparagraph:P039:end

% mdparagraph:P040:begin
\itempara{2}{\textbf{Hūlic is ðæt scip?} — What sort of ship is that?}
% mdparagraph:P040:end

% mdparagraph:P041:begin
\textbf{Target:} attributive \textbf{hwæðer} agrees with nominative masculine singular \textbf{mann}; the stated two-man context supplies “of the two.” This does not require an unassigned partitive phrase. Predicative \textbf{hūlic} is nominative neuter singular, agreeing with \textbf{ðæt scip}, and asks about its kind or character. General \textbf{hwilc} could ask a related question, but would miss the expressly requested word in either part.
% mdparagraph:P041:end

\end{minipage}\par\addvspace{4pt}

\par\begin{minipage}{\linewidth}\raggedright
% item: NEW-U08-Q05
% key: KEY-NEW-U08-Q05
% mdparagraph:P042:begin
\textbf{Q05. How and why.}
% mdparagraph:P042:end

% mdparagraph:P043:begin
\textbf{Hū swimð se mann?} — How does the man swim? \textbf{Hwȳ swimð se mann?} — Why does the man swim?
% mdparagraph:P043:end

% mdparagraph:P044:begin
\textbf{Target:} \textbf{hū} requests manner; \textbf{hwȳ} requests reason. Both are invariant question adverbs here, and \textbf{swimð} remains third-singular present. Rendering both questions “How?” loses the distinction. Lexical \textbf{hwȳ} “why” does not require the learner to construct an instrumental noun phrase.
% mdparagraph:P044:end

\end{minipage}\par\addvspace{4pt}

\par\begin{minipage}{\linewidth}\raggedright
% item: NEW-U08-Q06
% key: KEY-NEW-U08-Q06
% mdparagraph:P045:begin
\textbf{Q06. Any fish?}
% mdparagraph:P045:end

% mdparagraph:P046:begin
\textbf{Hæfst ðū ǣnigne fisc?} — Do you have any fish?
% mdparagraph:P046:end

% mdparagraph:P047:begin
\textbf{Target:} second-singular \textbf{hæfst} forms the question without an auxiliary; \textbf{ǣnigne fisc} is masculine accusative singular, as required. English “do” supports the English question, but Old English does not require a corresponding \textbf{dōn}. A plural-fish question could be meaningful Old English, but would miss the explicit singular \textbf{fisc} target here.
% mdparagraph:P047:end

\end{minipage}\par\addvspace{4pt}

\reserve{7\baselineskip}
% mdparagraph:P048:begin
\section*{N. Pronouns and agreement}
% mdparagraph:P048:end

\par\begin{minipage}{\linewidth}\raggedright
% item: NEW-U08-N01
% key: KEY-NEW-U08-N01
% mdparagraph:P049:begin
\textbf{N01. Some man; some men.}
% mdparagraph:P049:end

% mdparagraph:P050:begin
\textbf{Sum mann hæfð scip.} — A certain man has a ship. \textbf{Sume męnn habbað scipu.} — Some men have ships.
% mdparagraph:P050:end

% mdparagraph:P051:begin
\textbf{Target:} \textbf{sum mann} is nominative masculine singular; \textbf{sume męnn} is nominative masculine plural. The subjects govern third-singular \textbf{hæfð} and plural \textbf{habbað}. The objects are accusative neuter singular \textbf{scip} and plural \textbf{scipu}. “Some man” is also a defensible rendering of \textbf{sum mann}. The second sentence does not specify that each man owns exactly one ship; do not add that distribution. \textbf{Męnn} is also a dative singular form elsewhere, but its agreement here establishes nominative plural.
% mdparagraph:P051:end

\end{minipage}\par\addvspace{4pt}

\par\begin{minipage}{\linewidth}\raggedright
% item: NEW-U08-N02
% key: KEY-NEW-U08-N02
% mdparagraph:P052:begin
\textbf{N02. Each man; each companion.}
% mdparagraph:P052:end

% mdparagraph:P053:begin
\textbf{Gehwilc mann siehð scip.} — Each man sees a ship. \textbf{Gehwilc gefēra siehð scip.} — Each companion sees a ship.
% mdparagraph:P053:end

% mdparagraph:P054:begin
\textbf{Target:} singular \textbf{gehwilc} with a singular masculine noun, taking third-singular \textbf{siehð}. The weak noun \textbf{gefēra} has nominative singular \textbf{-a}, but this does not make its agreeing interrogative-derived indefinite \textbf{gehwilc} weak or plural. “Every man/\allowbreak{}companion” preserves the sense. A sentence using \textbf{ǣlc} could be grammatical, but would miss the requested \textbf{gehwilc}.
% mdparagraph:P054:end

\end{minipage}\par\addvspace{4pt}

\par\begin{minipage}{\linewidth}\raggedright
% item: NEW-U08-N03
% key: KEY-NEW-U08-N03
% mdparagraph:P055:begin
\textbf{N03. Other and another.}
% mdparagraph:P055:end

% mdparagraph:P056:begin
\textbf{Se ōðer mann hæfð ōðer scip.} — The other man has another ship.
% mdparagraph:P056:end

% mdparagraph:P057:begin
\textbf{Target:} first \textbf{ōðer} is nominative masculine singular with \textbf{mann}; second \textbf{ōðer} is accusative neuter singular with \textbf{scip}. The article does not give this adjective a new weak form. The sentence distinguishes a man and a ship from others already in view; it need not place either second in a numbered series. The earlier ordinal sense “second” is possible in a suitable context, but is not the meaning requested here.
% mdparagraph:P057:end

\end{minipage}\par\addvspace{4pt}

\par\begin{minipage}{\linewidth}\raggedright
% item: NEW-U08-N04
% key: KEY-NEW-U08-N04
% mdparagraph:P058:begin
\textbf{N04. The same man in two cases.}
% mdparagraph:P058:end

% mdparagraph:P059:begin
\textbf{Ic geseah ðone ilcan mann.} — I saw the same man. \textbf{Ic sęlle ðone fisc ðǣm ilcan męnn.} — I give the fish to the same man.
% mdparagraph:P059:end

% mdparagraph:P060:begin
\textbf{Target:} \textbf{ðone ilcan mann} is masculine accusative singular, the person seen. In the second sentence, \textbf{ðone fisc} is masculine accusative singular, the thing given, and \textbf{ðǣm ilcan męnn} is masculine dative singular, the recipient. \textbf{Ic} is nominative singular in each sentence. Weak \textbf{ilcan} is identical in these two oblique cases, but the article and noun differ. Do not retain accusative \textbf{ðone mann} for the recipient merely because English uses “the same man” in both clauses.
% mdparagraph:P060:end

\end{minipage}\par\addvspace{4pt}

\par\begin{minipage}{\linewidth}\raggedright
% item: NEW-U08-N05
% key: KEY-NEW-U08-N05
% mdparagraph:P061:begin
\textbf{N05. Subject emphasis and a reflexive object.}
% mdparagraph:P061:end

% mdparagraph:P062:begin
\textbf{Hē self geseah ðæt scip.} — He himself saw that ship. \textbf{Ic nęrie mē selfne.} — I save myself.
% mdparagraph:P062:end

% mdparagraph:P063:begin
\textbf{Target:} strong nominative masculine singular \textbf{self} emphasizes subject \textbf{hē} in the first sentence; the object is the ship. In the second, \textbf{mē} is an accusative object referring back to subject \textbf{ic}, reinforced by masculine accusative singular \textbf{selfne} because the specified speaker is male. \textbf{Mē} can have dative uses elsewhere, but \textbf{nęrian} takes the person saved as its accusative object. A simple reflexive \textbf{mē} could be valid language, but omitting \textbf{selfne} would miss the explicit target.
% mdparagraph:P063:end

\end{minipage}\par\addvspace{4pt}

\par\begin{minipage}{\linewidth}\raggedright
% item: NEW-U08-N06
% key: KEY-NEW-U08-N06
% mdparagraph:P064:begin
\textbf{N06. Such men.}
% mdparagraph:P064:end

% mdparagraph:P065:begin
\textbf{Swilce męnn sindon gōde.} — Such men are good.
% mdparagraph:P065:end

% mdparagraph:P066:begin
\textbf{Target:} \textbf{swilce} is nominative masculine plural of inflected \textbf{swilc}, agreeing with \textbf{męnn}. It means “such,” describing the men. The invariant adverb \textbf{swilce} “likewise” has the same spelling in this cell but does not agree with a noun. “The men are likewise good” would select a different function and lose the assigned noun-phrase meaning.
% mdparagraph:P066:end

\end{minipage}\par\addvspace{4pt}

\par\begin{minipage}{\linewidth}\raggedright
% item: NEW-U08-N07
% key: KEY-NEW-U08-N07
% mdparagraph:P067:begin
\textbf{N07. Four noun phrases.}
% mdparagraph:P067:end

\begin{itemize}
% mdparagraph:P068:begin
\item \textbf{Sumra manna} — of some men: masculine genitive plural; both words mark that case and number.
% mdparagraph:P068:end
% mdparagraph:P069:begin
\item \textbf{Ǣnigum gefēran} — to any companion: masculine dative singular; strong \textbf{ǣnigum} agrees with weak noun \textbf{gefēran}.
% mdparagraph:P069:end
% mdparagraph:P070:begin
\item \textbf{Ðæt ilce scip} — that same ship: neuter accusative singular; weak \textbf{ilce} agrees with \textbf{scip}, and \textbf{ðæt} marks the same features.
% mdparagraph:P070:end
% mdparagraph:P071:begin
\item \textbf{Swilce lȳtle fuglas} — such little birds: masculine accusative plural throughout; \textbf{lȳtle} is strong in this indefinite phrase.
% mdparagraph:P071:end
\end{itemize}

% mdparagraph:P072:begin
\textbf{Target:} use the requested cases directly. English “of/\allowbreak{}to” identifies a genitive or dative phrase here; it does not require adding a preposition. \textbf{Ðæt} could also be nominative in another clause, but the task assigns accusative. Weak plural \textbf{lȳtlan} belongs to a different adjective environment, such as Cook I.8.5 with \textbf{ðās}.
% mdparagraph:P072:end

\end{minipage}\par\addvspace{4pt}

\reserve{7\baselineskip}
% mdparagraph:P073:begin
\section*{L. Relative clauses}
% mdparagraph:P073:end

% mdparagraph:P074:begin
For L01–L08, distinguish the antecedent's role outside the relative clause from the relative's role within it. Invariant \textbf{ðe} has a syntactic role without showing a case ending. With an inflected relative, gender and number follow the antecedent, while the selected ordinary case follows the internal role. Models often put the finite verb late in the relative; this is a useful pattern, not an absolute rule governing every attested clause. [1]
% mdparagraph:P074:end

\par\begin{minipage}{\linewidth}\raggedright
% item: NEW-U08-L01
% key: KEY-NEW-U08-L01
% mdparagraph:P075:begin
\textbf{L01. A relative subject.}
% mdparagraph:P075:end

% mdparagraph:P076:begin
\textbf{Se mann, ðe ðæt scip hæfð, is gōd.} — The man who has that ship is good.
% mdparagraph:P076:end

% mdparagraph:P077:begin
\textbf{Target:} \textbf{se mann} is the main-clause subject and the antecedent of \textbf{ðe}. Within the relative, \textbf{ðe} supplies the subject of \textbf{hæfð}; \textbf{ðæt scip} is its accusative object. The owner is the man, not the ship. An inflected relative could express the same relation, but this task expressly asks for \textbf{ðe}.
% mdparagraph:P077:end

\end{minipage}\par\addvspace{4pt}

\par\begin{minipage}{\linewidth}\raggedright
% item: NEW-U08-L02
% key: KEY-NEW-U08-L02
% mdparagraph:P078:begin
\textbf{L02. A relative object.}
% mdparagraph:P078:end

% mdparagraph:P079:begin
\textbf{Ic geseah ðæt scip, ðe se mann hæfð.} — I saw the ship which the man has.
% mdparagraph:P079:end

% mdparagraph:P080:begin
\textbf{Target:} retain first-singular past \textbf{geseah} in the main clause and third-singular present \textbf{hæfð} in the relative. Antecedent \textbf{ðæt scip} is the thing seen; within the relative \textbf{ðe} is the object possessed, while \textbf{se mann} is subject. Making the man the main-clause subject or making the ship the subject of \textbf{hæfð} changes the participants and meaning.
% mdparagraph:P080:end

\end{minipage}\par\addvspace{4pt}

\par\begin{minipage}{\linewidth}\raggedright
% item: NEW-U08-L03
% key: KEY-NEW-U08-L03
% mdparagraph:P081:begin
\textbf{L03. Dative antecedent, nominative relative.}
% mdparagraph:P081:end

% mdparagraph:P082:begin
\textbf{Ic sęlle ðone fisc ðǣm męnn, se ðe mē healp.} — I give that fish to the man who helped me.
% mdparagraph:P082:end

% mdparagraph:P083:begin
\textbf{Target:} antecedent \textbf{ðǣm męnn} is dative singular as recipient of \textbf{sęlle}. Relative \textbf{se} is nominative masculine singular as subject of \textbf{healp}, followed by invariant \textbf{ðe}. The person helped, \textbf{mē}, is dative because \textbf{helpan} governs that case. An invariant \textbf{ðe} relative alone could be grammatical, but would omit the requested inflected demonstrative. Copying the antecedent's dative into the relative would miss this task's ordinary internal-case target.
% mdparagraph:P083:end

\end{minipage}\par\addvspace{4pt}

\par\begin{minipage}{\linewidth}\raggedright
% item: NEW-U08-L04
% key: KEY-NEW-U08-L04
% mdparagraph:P084:begin
\textbf{L04. Accusative antecedent, dative relative.}
% mdparagraph:P084:end

% mdparagraph:P085:begin
\textbf{Ic geseah ðone mann, ðǣm ðe ic healp.} — I saw the man whom I helped.
% mdparagraph:P085:end

% mdparagraph:P086:begin
\textbf{Target:} \textbf{ðone mann} is accusative singular as object of \textbf{geseah}. Relative \textbf{ðǣm} is dative masculine singular because the man is the person helped within the relative; \textbf{ic} is its nominative subject. English “whom” does not make every Old English personal relative accusative. An accusative \textbf{ðone} relative would fail the assigned ordinary government of \textbf{helpan}; the recognition exception does not remove that production target.
% mdparagraph:P086:end

\end{minipage}\par\addvspace{4pt}

\par\begin{minipage}{\linewidth}\raggedright
% item: NEW-U08-L05
% key: KEY-NEW-U08-L05
% mdparagraph:P087:begin
\textbf{L05. A singular addressee.}
% mdparagraph:P087:end

% mdparagraph:P088:begin
\textbf{Ðū, ðe mē hulpe, sęle mē ðone fisc.} — You, who helped me, give me that fish.
% mdparagraph:P088:end

% mdparagraph:P089:begin
\textbf{Target:} \textbf{hulpe} is second-singular preterite indicative, agreeing with personal antecedent \textbf{ðū}; \textbf{sęle} is a second-singular command to the same person. The first \textbf{mē} is the dative person helped, the second the dative recipient of the fish; \textbf{ðone fisc} is accusative singular. Third-singular \textbf{healp} would lose the selected second-person agreement. The relative “who” does not require interrogative \textbf{hwā}.
% mdparagraph:P089:end

\end{minipage}\par\addvspace{4pt}

\par\begin{minipage}{\linewidth}\raggedright
% item: NEW-U08-L06
% key: KEY-NEW-U08-L06
% mdparagraph:P090:begin
\textbf{L06. Two or more addressees.}
% mdparagraph:P090:end

% mdparagraph:P091:begin
\textbf{Gē, ðe mē hulpon, sęllað mē ðone fisc.} — You, who helped me, give me that fish.
% mdparagraph:P091:end

% mdparagraph:P092:begin
\textbf{Target:} replace \textbf{ðū → gē}, \textbf{hulpe → hulpon}, and \textbf{sęle → sęllað}. The plural forms address two or more people. Invariant \textbf{ðe} and first-singular \textbf{mē} remain unchanged, and \textbf{ðone fisc} stays singular as requested. A specifically dual pronoun could identify exactly two addressees where that is intended and would still take plural finite forms; it does not justify retaining singular \textbf{hulpe/\allowbreak{}sęle}.
% mdparagraph:P092:end

\end{minipage}\par\addvspace{4pt}

\par\begin{minipage}{\linewidth}\raggedright
% item: NEW-U08-L07
% key: KEY-NEW-U08-L07
% mdparagraph:P093:begin
\textbf{L07. Subject, object, and possessor within a relative.}
% mdparagraph:P093:end

% mdparagraph:P094:begin
In \textbf{Se mann, ðe ðæt scip hæfð, is mīn gefēra}, the antecedent is nominative singular \textbf{se mann}, subject of \textbf{is}. Relative \textbf{ðe} is subject of \textbf{hæfð}; \textbf{ðæt scip} is its object. The meaning is “The man who has that ship is my companion.”
% mdparagraph:P094:end

% mdparagraph:P095:begin
In \textbf{Ic geseah ðæt scip, ðe se mann hæfð}, the antecedent is accusative singular \textbf{ðæt scip}, object of \textbf{geseah}. Relative \textbf{ðe} is object of \textbf{hæfð}; \textbf{se mann} is its subject. The meaning is “I saw the ship which the man has.” \textbf{Ðe} does not change its spelling to mark either internal role.
% mdparagraph:P095:end

% mdparagraph:P096:begin
\textbf{Se mann, ðæs scip ic geseah, is gōd.} — The man whose ship I saw is good. Antecedent \textbf{se mann} is nominative masculine singular, the main-clause subject. Relative \textbf{ðæs} is genitive masculine singular: it refers to the man as possessor of \textbf{scip}. \textbf{Scip} is accusative neuter singular, the thing seen by subject \textbf{ic}. Thus \textbf{ðæs} and \textbf{scip} need not share a case; possession within an object phrase requires genitive \textbf{ðæs}, not an accusative relative copied from the object noun.
% mdparagraph:P096:end

% mdparagraph:P097:begin
\textbf{Target:} identify the internal function separately in all three clauses. The genitive example does not involve case attraction: the man is nominative outside the relative, while the relative is genitive because it marks possession within it.
% mdparagraph:P097:end

\end{minipage}\par\addvspace{4pt}

\par\begin{minipage}{\linewidth}\raggedright
% item: NEW-U08-L08
% key: KEY-NEW-U08-L08
% mdparagraph:P098:begin
\textbf{L08. Correcting an overbroad rule.}
% mdparagraph:P098:end

% mdparagraph:P099:begin
An inflected relative does \textbf{not always} have its antecedent's case. The ordinary production rule is to choose case from its role within the relative clause. In L03, dative recipient \textbf{ðǣm męnn} is followed by nominative relative subject \textbf{se}. In L04, accusative object \textbf{ðone mann} is followed by dative \textbf{ðǣm}, governed by \textbf{helpan} in the relative.
% mdparagraph:P099:end

% mdparagraph:P100:begin
\textbf{Target:} retain agreement in gender and number while separating those two case decisions. Attested \textbf{case attraction} sometimes gives a relative the antecedent's case instead; the lesson presents this as a recognition exception. It neither makes the learner's “always” rule correct nor licenses ignoring the explicit ordinary-case target of L03 and L04. Conversely, claiming that case attraction is impossible would overstate the ordinary rule.
% mdparagraph:P100:end

\end{minipage}\par\addvspace{4pt}

\reserve{7\baselineskip}
% mdparagraph:P101:begin
\section*{C. Lexical constructions}
% mdparagraph:P101:end

\par\begin{minipage}{\linewidth}\raggedright
% item: NEW-U08-C01
% key: KEY-NEW-U08-C01
% mdparagraph:P102:begin
\textbf{C01. Finish and fill.}
% mdparagraph:P102:end

% mdparagraph:P103:begin
\textbf{Hīe gefyllað ðæt weorc.} — They finish the work. \textbf{Hīe gefyllað ðā wæteru.} — They fill the waters.
% mdparagraph:P103:end

% mdparagraph:P104:begin
\textbf{Target:} the same plural present form \textbf{gefyllað}, with distinct lexical senses selected by the objects: finish/\allowbreak{}accomplish the work; fill the waters. \textbf{Ðæt weorc} is accusative neuter singular, \textbf{ðā wæteru} accusative neuter plural. Familiarity with the first sense does not erase the second. Retain the course's selected plural \textbf{wæteru} rather than silently making “the waters” singular.
% mdparagraph:P104:end

\end{minipage}\par\addvspace{4pt}

\par\begin{minipage}{\linewidth}\raggedright
% item: NEW-U08-C02
% key: KEY-NEW-U08-C02
% mdparagraph:P105:begin
\textbf{C02. Having in one's control.}
% mdparagraph:P105:end

% mdparagraph:P106:begin
\textbf{Se mann hæfð eall nīetenu on his gewealde.} — The man has all the animals in his control. \textbf{Hafa eall nīetenu on ðīnum gewealde.} — Have all the animals in your control.
% mdparagraph:P106:end

% mdparagraph:P107:begin
\textbf{Target:} \textbf{eall nīetenu} is the accusative neuter plural entity controlled. The idiom uses \textbf{on + dative gewealde}. Third-person possessor \textbf{his} is invariant, while \textbf{ðīnum} agrees in dative neuter singular with \textbf{gewealde}. \textbf{Hafa} is the command to one person; the plural command would miss that stated addressee. Previously taught \textbf{dēor} can also express “animals” with the appropriate neuter plural agreement. A literal spatial rendering with “under” does not supply this control idiom.
% mdparagraph:P107:end

\end{minipage}\par\addvspace{4pt}

\par\begin{minipage}{\linewidth}\raggedright
% item: NEW-U08-C03
% key: KEY-NEW-U08-C03
% mdparagraph:P108:begin
\textbf{C03. Living one's life.}
% mdparagraph:P108:end

% mdparagraph:P109:begin
\textbf{Ic libbe mīn līf.} — I live my life. \textbf{Wē libbað ūru līf.} — We live our lives.
% mdparagraph:P109:end

% mdparagraph:P110:begin
\textbf{Target:} \textbf{mīn līf} is accusative neuter singular; \textbf{ūru līf} is accusative neuter plural. \textbf{Libban} here takes a cognate accusative, an object expressing the life lived. The individual-life sense permits the plural. First-singular \textbf{libbe} changes to first-plural \textbf{libbað}, and the selected agreeing possessive \textbf{ūru} confirms plural \textbf{līf}, despite its unchanged noun ending. \textbf{Ūre} is not the selected neuter plural possessive cell for this task.
% mdparagraph:P110:end

\end{minipage}\par\addvspace{4pt}

\par\begin{minipage}{\linewidth}\raggedright
% item: NEW-U08-C04
% key: KEY-NEW-U08-C04
% mdparagraph:P111:begin
\textbf{C04. Appears; let it appear.}
% mdparagraph:P111:end

% mdparagraph:P112:begin
\textbf{Drȳgnes ætēowað.} — Dry land appears. \textbf{Ætēowie drȳgnes.} — Let dry land appear.
% mdparagraph:P112:end

% mdparagraph:P113:begin
\textbf{Target:} third-singular present indicative \textbf{ætēowað} states the event; third-singular present subjunctive \textbf{ætēowie} expresses the jussive. The dry land is subject, not an object being displayed by someone else. Although \textbf{ætēowie} also spells a first-singular indicative form, its third-singular subject and jussive force determine its use here. Do not replace “let” with an unrelated permission verb or replace intransitive “appear” with transitive “reveal.”
% mdparagraph:P113:end

\end{minipage}\par\addvspace{4pt}

\reserve{7\baselineskip}
% mdparagraph:P114:begin
\section*{Cook I.8}
% mdparagraph:P114:end

% mdparagraph:P115:begin
The English prompts below retain Cook's complete set and numbering. Where English leaves the addressee's number or the number of “fish” open, the model makes one choice; a defensible alternative must carry that choice through its agreement. An expressly singular command or supplied singular form still fixes the number addressed. [2]
% mdparagraph:P115:end

\par\begin{minipage}{\linewidth}\raggedright
% item: EX-I-08-01
% key: KEY-EX-I-08-01
% mdparagraph:P116:begin
\textbf{I.8.1. Your works are good.}
% mdparagraph:P116:end

% mdparagraph:P117:begin
\textbf{Ēowru weorc sindon gōd.}
% mdparagraph:P117:end

% mdparagraph:P118:begin
\textbf{Target:} nominative neuter plural \textbf{ēowru weorc}, plural \textbf{sindon}, and strong predicate neuter plural \textbf{gōd}. The model addresses several people; \textbf{Ðīn weorc sindon gōd} is valid for one possessor because Cook's “your” does not decide that number. The number of possessors does not make the works singular. Selected neuter plural \textbf{ēowru} differs from \textbf{ēowre}, and zero-ending predicate \textbf{gōd} differs from masculine plural \textbf{gōde}.
% mdparagraph:P118:end

\end{minipage}\par\addvspace{4pt}

\par\begin{minipage}{\linewidth}\raggedright
% item: EX-I-08-02
% key: KEY-EX-I-08-02
% mdparagraph:P119:begin
\textbf{I.8.2. His tree is fruit-bearing.}
% mdparagraph:P119:end

% mdparagraph:P120:begin
\textbf{His trēo is æppelbǣre.}
% mdparagraph:P120:end

% mdparagraph:P121:begin
\textbf{Target:} invariant possessor \textbf{his}, neuter singular \textbf{trēo}, singular \textbf{is}, and predicate adjective \textbf{æppelbǣre}. The essential meaning is that his tree bears fruit. Selected \textbf{trēo} and documented \textbf{trēow} are spelling alternatives, not different numbers; use the selected form when reproducing the course paradigm. Do not add a possessive agreement ending to \textbf{his}.
% mdparagraph:P121:end

\end{minipage}\par\addvspace{4pt}

\par\begin{minipage}{\linewidth}\raggedright
% item: EX-I-08-03
% key: KEY-EX-I-08-03
% mdparagraph:P122:begin
\textbf{I.8.3. Thou art the man.}
% mdparagraph:P122:end

% mdparagraph:P123:begin
\textbf{Ðū eart se mann.}
% mdparagraph:P123:end

% mdparagraph:P124:begin
\textbf{Target:} second-singular \textbf{ðū eart} and nominative predicate \textbf{se mann}. This identifies the addressee as the man. Accusative \textbf{ðone mann} would treat the predicate as though “be” took a direct object; third-singular \textbf{is} would lose the stated person.
% mdparagraph:P124:end

\end{minipage}\par\addvspace{4pt}

\par\begin{minipage}{\linewidth}\raggedright
% item: EX-I-08-04
% key: KEY-EX-I-08-04
% mdparagraph:P125:begin
\textbf{I.8.4. This year is the fourth.}
% mdparagraph:P125:end

% mdparagraph:P126:begin
\textbf{Ðis gēar is ðæt fēorðe.}
% mdparagraph:P126:end

% mdparagraph:P127:begin
\textbf{Target:} nominative neuter singular \textbf{ðis gēar} and the definite ordinal \textbf{ðæt fēorðe}, with weak neuter singular \textbf{fēorðe}. “Fourth” identifies a place in a sequence; it is not the cardinal “four.” Masculine \textbf{se} would not agree with neuter \textbf{gēar}.
% mdparagraph:P127:end

\end{minipage}\par\addvspace{4pt}

\par\begin{minipage}{\linewidth}\raggedright
% item: EX-I-08-05
% key: KEY-EX-I-08-05
% mdparagraph:P128:begin
\textbf{I.8.5. God made these little birds.}
% mdparagraph:P128:end

% mdparagraph:P129:begin
\textbf{God geworhte ðās lȳtlan fuglas.}
% mdparagraph:P129:end

% mdparagraph:P130:begin
\textbf{Target:} third-singular preterite \textbf{geworhte} and accusative masculine plural \textbf{ðās lȳtlan fuglas}. The demonstrative gives the ordinary descriptive adjective its weak form \textbf{lȳtlan}. The essential meaning is God's making these particular small birds. \textbf{Gescēop} can also express creation here; it must retain the singular past subject and plural object. Strong \textbf{lȳtle} is the selected form in an indefinite phrase such as “such little birds,” not this definite environment.
% mdparagraph:P130:end

\end{minipage}\par\addvspace{4pt}

\par\begin{minipage}{\linewidth}\raggedright
% item: EX-I-08-06
% key: KEY-EX-I-08-06
% mdparagraph:P131:begin
\textbf{I.8.6. He created me likewise.}
% mdparagraph:P131:end

% mdparagraph:P132:begin
\textbf{Hē gescēop mē ēac swilce.}
% mdparagraph:P132:end

% mdparagraph:P133:begin
\textbf{Target:} third-singular preterite \textbf{gescēop}, accusative object \textbf{mē}, and the adverbial “also/\allowbreak{}likewise” expression \textbf{ēac swilce}. The speaker is the person created. Invariant adverb \textbf{swilce} here is not an inflected adjective modifying \textbf{mē}. A valid rendering with the established “likewise” adverb alone need not be rejected for omitting the reinforcing “also.”
% mdparagraph:P133:end

\end{minipage}\par\addvspace{4pt}

\par\begin{minipage}{\linewidth}\raggedright
% item: EX-I-08-07
% key: KEY-EX-I-08-07
% mdparagraph:P134:begin
\textbf{I.8.7. Our cattle (plur.) see the grass.}
% mdparagraph:P134:end

% mdparagraph:P135:begin
\textbf{Ūru nīetenu sēoð ðæt gærs.}
% mdparagraph:P135:end

% mdparagraph:P136:begin
\textbf{Target:} nominative neuter plural \textbf{ūru nīetenu}, plural present indicative \textbf{sēoð}, and accusative neuter singular \textbf{ðæt gærs}. The cattle, not the grass, do the seeing. The selected possessive is \textbf{ūru} in this neuter plural cell. \textbf{Sēon} is an infinitive or a subjunctive form in other contexts; it does not supply the present indicative required by this statement.
% mdparagraph:P136:end

\end{minipage}\par\addvspace{4pt}

\par\begin{minipage}{\linewidth}\raggedright
% item: EX-I-08-08
% key: KEY-EX-I-08-08
% mdparagraph:P137:begin
\textbf{I.8.8. We live our lives (168.1).}
% mdparagraph:P137:end

% mdparagraph:P138:begin
\textbf{Wē libbað ūru līf.}
% mdparagraph:P138:end

% mdparagraph:P139:begin
\textbf{Target:} first-plural present \textbf{libbað} with the cognate accusative \textbf{ūru līf}, accusative neuter plural. The several people live their respective lives; \textbf{līf} has the assigned individual-life sense. Its zero plural ending does not make it singular. \textbf{Ūru}, rather than an assumed \textbf{ūre}, is the selected neuter plural possessive. Cook's reference is \textbf{168.1}, the cognate-accusative discussion.
% mdparagraph:P139:end

\end{minipage}\par\addvspace{4pt}

\par\begin{minipage}{\linewidth}\raggedright
% item: EX-I-08-09
% key: KEY-EX-I-08-09
% mdparagraph:P140:begin
\textbf{I.8.9. His seeds grow in the earth.}
% mdparagraph:P140:end

% mdparagraph:P141:begin
\textbf{His sǣd weaxað on ðǣre eorðan.}
% mdparagraph:P141:end

% mdparagraph:P142:begin
\textbf{Target:} nominative neuter plural \textbf{sǣd}, plural present \textbf{weaxað}, and dative singular \textbf{ðǣre eorðan} after locative \textbf{on}. \textbf{His} is invariant. The unchanged noun ending in \textbf{sǣd} does not override the plural meaning and verb; third-singular \textbf{wiext} would describe one seed instead.
% mdparagraph:P142:end

\end{minipage}\par\addvspace{4pt}

\par\begin{minipage}{\linewidth}\raggedright
% item: EX-I-08-10
% key: KEY-EX-I-08-10
% mdparagraph:P143:begin
\textbf{I.8.10. The animals rest (themselves).}
% mdparagraph:P143:end

% mdparagraph:P144:begin
\textbf{Ðā nīetenu geręstað hīe.}
% mdparagraph:P144:end

% mdparagraph:P145:begin
\textbf{Target:} plural subject \textbf{ðā nīetenu}, plural present \textbf{geręstað}, and accusative reflexive \textbf{hīe}, referring to those same animals. Previously learned \textbf{dēor} is also a valid neuter plural noun for “animals” here. The parenthetic “themselves” points to the established reflexive construction; \textbf{hīe} is not a new group of animals receiving an action from the first. Emphatic \textbf{self} is not required merely because English uses a reflexive pronoun.
% mdparagraph:P145:end

\end{minipage}\par\addvspace{4pt}

\par\begin{minipage}{\linewidth}\raggedright
% item: EX-I-08-11
% key: KEY-EX-I-08-11
% mdparagraph:P146:begin
\textbf{I.8.11. Fill the waters.}
% mdparagraph:P146:end

% mdparagraph:P147:begin
\textbf{Gefyllað ðā wæteru.}
% mdparagraph:P147:end

% mdparagraph:P148:begin
\textbf{Target:} the command of \textbf{gefyllan} in its “fill” sense, with accusative neuter plural \textbf{ðā wæteru}. The model addresses several people. \textbf{Gefyl ðā wæteru} is equally defensible for one addressee, since the English prompt does not fix the command's number. Do not confuse the number commanded with the number of waters, or substitute the “finish the work” sense.
% mdparagraph:P148:end

\end{minipage}\par\addvspace{4pt}

\par\begin{minipage}{\linewidth}\raggedright
% item: EX-I-08-12
% key: KEY-EX-I-08-12
% mdparagraph:P149:begin
\textbf{I.8.12. Let dry land appear.}
% mdparagraph:P149:end

% mdparagraph:P150:begin
\textbf{Ætēowie drȳgnes.}
% mdparagraph:P150:end

% mdparagraph:P151:begin
\textbf{Target:} present subjunctive \textbf{ætēowie} with third-singular subject \textbf{drȳgnes}, giving jussive “let … appear.” The statement \textbf{Drȳgnes ætēowað} means “Dry land appears” and changes the force. The source and dictionary identify the intransitive weak-II verb \textbf{ætēowian}; the transitive “reveal, display” verb does not express this same subject relation. A different defensible order must preserve that mood and relation.
% mdparagraph:P151:end

\end{minipage}\par\addvspace{4pt}

\par\begin{minipage}{\linewidth}\raggedright
% item: EX-I-08-13
% key: KEY-EX-I-08-13
% mdparagraph:P152:begin
\textbf{I.8.13. Place (gesęte) the fish in the sea.}
% mdparagraph:P152:end

% mdparagraph:P153:begin
\textbf{Gesęte ðone fisc on ðǣm sǣ.}
% mdparagraph:P153:end

% mdparagraph:P154:begin
\textbf{Target:} use the supplied singular command \textbf{gesęte}, the accusative thing placed, and the locally supplied \textbf{on + dative} placement frame. The model takes “fish” as singular; \textbf{Gesęte ðā fiscas on ðǣm sǣ} is valid for plural fish, also allowed by the English wording. The supplied command still addresses one person. The attested comparable placement uses a dative location; the fact that someone places an object does not by itself force accusative after \textbf{on}. Neither the prefixed command nor the whole verb system is an unaided-recall requirement.
% mdparagraph:P154:end

\end{minipage}\par\addvspace{4pt}

\par\begin{minipage}{\linewidth}\raggedright
% item: EX-I-08-14
% key: KEY-EX-I-08-14
% mdparagraph:P155:begin
\textbf{I.8.14. Have (sing.) all spirits under [in] your control.}
% mdparagraph:P155:end

% mdparagraph:P156:begin
\textbf{Hafa ealle gāstas on ðīnum gewealde.}
% mdparagraph:P156:end

% mdparagraph:P157:begin
\textbf{Target:} singular command \textbf{hafa}; accusative masculine plural \textbf{ealle gāstas}, the spirits controlled; \textbf{on ðīnum gewealde}, with dative neuter singular possessive and noun. The bracketed “in” points to the control expression, not to a separate physical location below the addressee. Plural \textbf{habbað} would miss the explicitly singular command, and masculine plural \textbf{ealle} must not be copied from the neuter plural \textbf{eall} used with animals in C02.
% mdparagraph:P157:end

\end{minipage}\par\addvspace{4pt}

\par\begin{minipage}{\linewidth}\raggedright
% item: EX-I-08-15
% key: KEY-EX-I-08-15
% mdparagraph:P158:begin
\textbf{I.8.15. Let the stars shine in the heavens.}
% mdparagraph:P158:end

% mdparagraph:P159:begin
\textbf{Scīnen ðā steorran on ðǣm heofonum.}
% mdparagraph:P159:end

% mdparagraph:P160:begin
\textbf{Target:} present subjunctive plural \textbf{scīnen}, with the stars as third-plural subject, and dative plural location \textbf{on ðǣm heofonum}. This is a jussive concerning the stars, not a command addressed to an expressed second person. Indicative \textbf{scīnað} would state that they shine; preterite subjunctive \textbf{scinen}, with short \textbf{i}, would change the tense. Retain the present long vowel and subjunctive plural \textbf{-en}.
% mdparagraph:P160:end

\end{minipage}\par\addvspace{4pt}

\par\begin{minipage}{\linewidth}\raggedright
% item: EX-I-08-16
% key: KEY-EX-I-08-16
% mdparagraph:P161:begin
\textbf{I.8.16. I say indeed, Bless all men.}
% mdparagraph:P161:end

% mdparagraph:P162:begin
\textbf{Ic cweðe sōðlice: “Geblētsiað ealle męnn.”}
% mdparagraph:P162:end

% mdparagraph:P163:begin
\textbf{Target:} first-singular present “I say,” the adverb \textbf{sōðlice} “indeed,” and a direct quoted command with accusative masculine plural \textbf{ealle męnn}. The model addresses several people within the quotation. \textbf{Geblētsa ealle męnn} is a valid command to one person; the English does not specify that number. Previously taught \textbf{sęcge} is also a valid reporting verb. Keep the distinction between the speaker of \textbf{ic cweðe} and the people addressed by “Bless”; the infinitive alone would not supply the command.
% mdparagraph:P163:end

\end{minipage}\par\addvspace{4pt}

\par\begin{minipage}{\linewidth}\raggedright
% item: EX-I-08-17
% key: KEY-EX-I-08-17
% mdparagraph:P164:begin
\textbf{I.8.17. Thou, who madest life, bless me.}
% mdparagraph:P164:end

% mdparagraph:P165:begin
\textbf{Ðū, ðe līf gescēope, geblētsa mē.}
% mdparagraph:P165:end

% mdparagraph:P166:begin
\textbf{Target:} second-singular antecedent \textbf{ðū}, relative subject \textbf{ðe}, second-singular preterite \textbf{gescēope}, and singular command \textbf{geblētsa}. \textbf{Līf} is the accusative object of creation; \textbf{mē} is the accusative person to be blessed. \textbf{Geworhtest} “madest” is also defensible in the relative if the person and meaning are preserved. Third-singular \textbf{gescēop} would lose the selected agreement with \textbf{ðū}, and plural \textbf{geblētsiað} would contradict the explicitly singular “Thou.”
% mdparagraph:P166:end

\end{minipage}\par\addvspace{4pt}

\reserve{7\baselineskip}
% mdparagraph:P167:begin
\section*{D. The ship and the fish}
% mdparagraph:P167:end

\par\begin{minipage}{\linewidth}\raggedright
% item: NEW-U08-D01
% key: KEY-NEW-U08-D01
% mdparagraph:P168:begin
\textbf{D01. Complete translation.}
% mdparagraph:P168:end

% mdparagraph:P169:begin
Three men and one ship are on the shore. A certain man says, “Who is the man who has that ship?” His companion says, “The man is my companion. He has the fish which we see.”
% mdparagraph:P169:end

% mdparagraph:P170:begin
\textbf{First man:} “To whom do you give that fish?”
% mdparagraph:P170:end

% mdparagraph:P171:begin
\textbf{His companion:} “I give it to the man who helped me.”
% mdparagraph:P171:end

% mdparagraph:P172:begin
\textbf{First man:} “Which ship does he have?”
% mdparagraph:P172:end

% mdparagraph:P173:begin
\textbf{His companion:} “He has the little ship which you saw.”
% mdparagraph:P173:end

% mdparagraph:P174:begin
\textbf{Target:} preserve the opening scene and all six spoken turns. The first speaker asks his companion about the third man, the ship's possessor; the answering companion identifies that man as his own companion. \textbf{Wē} includes the speaker in “we see,” and \textbf{ðū} in the final relative addresses the first man. The fish are plural in \textbf{ðā fiscas}, but the later \textbf{ðone fisc/\allowbreak{}hine} selects one fish from that group; English “the fish” can conceal this number change. The recipient is identified as the man who helped the answering speaker. The account gives no names or further reason for the transfer, so do not invent them.
% mdparagraph:P174:end

% mdparagraph:P175:begin
“A man” or “some man” can render \textbf{sum mann}. English “that” or “the” may render the demonstrative/\allowbreak{}article according to the context; either must preserve the references. The final \textbf{gesāwe} is past “saw,” not present “see.”
% mdparagraph:P175:end

\end{minipage}\par\addvspace{4pt}

\par\begin{minipage}{\linewidth}\raggedright
% item: NEW-U08-D02
% key: KEY-NEW-U08-D02
% mdparagraph:P176:begin
\textbf{D02. The three direct questions.}
% mdparagraph:P176:end

% mdparagraph:P177:begin
\itempara{1}{\textbf{Hwā is se mann, ðe ðæt scip hæfð?} — Who is the man who has that ship? \textbf{Hwā} is nominative in a copular identification question: it asks the identity of the man described by \textbf{se mann} and its relative. It can be described as the interrogative predicate nominative; it is not an accusative object of \textbf{is}.}
% mdparagraph:P177:end

% mdparagraph:P178:begin
\itempara{2}{\textbf{Hwǣm sęlst ðū ðone fisc?} — To whom do you give that fish? \textbf{Hwǣm} is dative, asking for the recipient of \textbf{sęlst}; \textbf{ðone fisc} is the object given.}
% mdparagraph:P178:end

% mdparagraph:P179:begin
\itempara{3}{\textbf{Hwilc scip hæfð hē?} — Which ship does he have? \textbf{Hwilc} is accusative neuter singular, agreeing with object \textbf{scip}; it asks which ship is possessed. Subject \textbf{hē} is nominative.}
% mdparagraph:P179:end

% mdparagraph:P180:begin
\textbf{Target:} distinguish identification, recipient and object selection. A question word's position at the beginning does not make it nominative; the third question's \textbf{hwilc} happens to have the same nominative and accusative neuter spelling. The \textbf{ðe} inside the first question introduces a relative, not a fourth direct question.
% mdparagraph:P180:end

\end{minipage}\par\addvspace{4pt}

\par\begin{minipage}{\linewidth}\raggedright
% item: NEW-U08-D03
% key: KEY-NEW-U08-D03
% mdparagraph:P181:begin
\textbf{D03. The answer beginning Ic sęlle.}
% mdparagraph:P181:end

% mdparagraph:P182:begin
\textbf{Ic sęlle hine ðǣm męnn, se ðe mē healp} means “I give it to the man who helped me.”
% mdparagraph:P182:end

\begin{itemize}
% mdparagraph:P183:begin
\item \textbf{Hine:} accusative masculine singular, the fish being given; masculine agreement follows \textbf{fisc}, although English naturally uses “it.”
% mdparagraph:P183:end
% mdparagraph:P184:begin
\item \textbf{Ðǣm męnn:} dative masculine singular, the recipient in the main clause.
% mdparagraph:P184:end
% mdparagraph:P185:begin
\item \textbf{Se:} nominative masculine singular, the subject of \textbf{healp} inside the relative; it refers to that same man.
% mdparagraph:P185:end
% mdparagraph:P186:begin
\item \textbf{Mē:} dative singular, the person helped, governed by \textbf{helpan}; it refers to the answering speaker.
% mdparagraph:P186:end
\end{itemize}

% mdparagraph:P187:begin
\textbf{Target:} the relative identifies the recipient by his previous help to the speaker. He receives the fish and previously gave help; the speaker gives the fish and previously received help. Making \textbf{se} dative solely because \textbf{ðǣm męnn} is dative would miss the selected internal-case pattern. Translating \textbf{hine} as another man here would break the fish reference.
% mdparagraph:P187:end

\end{minipage}\par\addvspace{4pt}

\par\begin{minipage}{\linewidth}\raggedright
% item: NEW-U08-D04
% key: KEY-NEW-U08-D04
% mdparagraph:P188:begin
\textbf{D04. Two little ships.}
% mdparagraph:P188:end

% mdparagraph:P189:begin
\textbf{Hē hæfð twā lȳtlu scipu, ðā ðe ðū gesāwe.} — He has two little ships which you saw.
% mdparagraph:P189:end

% mdparagraph:P190:begin
\textbf{Target:} \textbf{twā}, \textbf{lȳtlu} and \textbf{scipu} are accusative neuter plural; \textbf{lȳtlu} is strong in the phrase headed by the numeral without the earlier article. Inflected relative \textbf{ðā} is likewise accusative neuter plural, the object of \textbf{gesāwe}, followed by unchanged \textbf{ðe}. Main subject \textbf{hē} still takes \textbf{hæfð}, while relative subject \textbf{ðū} still takes second-singular \textbf{gesāwe}. Do not pluralize either finite verb merely because the ships are plural. Retaining an explicitly definite phrase could support a different adjective environment, but agreement and the requested two ships must remain consistent.
% mdparagraph:P190:end

\end{minipage}\par\addvspace{4pt}

\reserve{7\baselineskip}
% mdparagraph:P191:begin
\section*{R. The merchant and his merchandise}
% mdparagraph:P191:end

% mdparagraph:P192:begin
These answers use the local help supplied with the exact four-turn excerpt. The translation includes the full selected passage; the following commodity catalogue is outside it. The nautical vocabulary, content clause, passive and closing participial construction remain supported reading material. [3]
% mdparagraph:P192:end

\par\begin{minipage}{\linewidth}\raggedright
% item: NEW-U08-R01
% key: KEY-NEW-U08-R01
% mdparagraph:P193:begin
\textbf{R01. Complete translation.}
% mdparagraph:P193:end

% mdparagraph:P194:begin
\textbf{Teacher:} “What do you say, merchant?”
% mdparagraph:P194:end

% mdparagraph:P195:begin
\textbf{Merchant:} “I say that I am useful both to the king and to nobles, and to wealthy people, and to all the people.”
% mdparagraph:P195:end

% mdparagraph:P196:begin
\textbf{Teacher:} “And how [are you useful]?”
% mdparagraph:P196:end

% mdparagraph:P197:begin
\textbf{Merchant:} “I board my ship with my cargo, and row across stretches of sea, and sell my goods, and buy valuable things which are not produced in this country. And I bring the goods here to you at great risk across the sea, and sometimes I suffer shipwreck, with the loss of all my things, barely escaping alive.”
% mdparagraph:P197:end

% mdparagraph:P198:begin
\textbf{Target:} preserve all four speakers' turns and the full sequence of claims: usefulness to the named groups; boarding with cargo; sea travel; selling; buying valuable things not produced locally; delivery here to the addressees; risk; occasional shipwreck and loss; and escape while alive. “The wealthy,” “the whole people,” “this land,” “merchandise” and “with difficulty escaping alive” are defensible English alternatives. The words restored in brackets explain the teacher's supplied elliptical question; they are not additional words in the Old English.
% mdparagraph:P198:end

% mdparagraph:P199:begin
\textbf{Cīepe} means “sell” here and \textbf{bycge} “buy”; reversing them changes the trade. Actual singular \textbf{hit} collectively refers to the goods, so English “it,” “the goods” or contextual “them” can preserve the reference without altering the source. \textbf{Ēow} is plural “you”; \textbf{hider} means “hither/\allowbreak{}here” as destination. The present forms describe the merchant's activity, not a past voyage invented by the translator. The closing words say that he barely escapes alive, not that he dies or escapes easily.
% mdparagraph:P199:end

\end{minipage}\par\addvspace{4pt}

\par\begin{minipage}{\linewidth}\raggedright
% item: NEW-U08-R02
% key: KEY-NEW-U08-R02
% mdparagraph:P200:begin
\textbf{R02. The teacher's questions.}
% mdparagraph:P200:end

% mdparagraph:P201:begin
\textbf{Hwæt sægst ðū, mangere?} — What do you say, merchant? \textbf{Hwæt} is the accusative thing said, the object of second-singular \textbf{sægst}; \textbf{ðū} is subject. Supplied \textbf{mangere} is the merchant's nominative form used in direct address, not another subject or the thing said.
% mdparagraph:P201:end

% mdparagraph:P202:begin
\textbf{And hū?} — And how? The local help restores “How [are you useful]?” from the preceding answer. Invariant \textbf{hū} asks how the merchant supplies the usefulness he has just claimed; his account of trade answers it.
% mdparagraph:P202:end

% mdparagraph:P203:begin
\textbf{Target:} identify these two direct questions and their different requested information. The content \textbf{ðæt} in \textbf{Ic sęcge ðæt …} introduces what the merchant says; it does not turn his answer into another question. No general independent construction of address or elliptical questions is demanded beyond the supplied forms and explanation.
% mdparagraph:P203:end

\end{minipage}\par\addvspace{4pt}

\par\begin{minipage}{\linewidth}\raggedright
% item: NEW-U08-R03
% key: KEY-NEW-U08-R03
% mdparagraph:P204:begin
\textbf{R03. The relative about the goods.}
% mdparagraph:P204:end

% mdparagraph:P205:begin
\textbf{Ðā on ðissum lande ne bēoð ācęnnede} means “which are not produced in this country.” Its antecedent is \textbf{ðing dēorwierðu}, “valuable things,” accusative neuter plural as the object of \textbf{bycge}. Relative \textbf{ðā} is nominative neuter plural, the subject of the passive within its clause. \textbf{Bēoð} is third-plural present, agreeing with those things; supplied \textbf{ācęnnede} is the plural past participle “produced.”
% mdparagraph:P205:end

% mdparagraph:P206:begin
\textbf{Target:} gender and number link relative and antecedent, but their cases differ because their clause roles differ. The goods are bought in the main clause and are not produced locally in the relative. Do not treat \textbf{ðā} as an accusative object merely because its antecedent is accusative, or omit \textbf{ne} and say that the goods are produced locally. The exact negative passive is supplied; no full productive system of \textbf{ācęnnan} is being tested.
% mdparagraph:P206:end

\end{minipage}\par\addvspace{4pt}

\par\begin{minipage}{\linewidth}\raggedright
% item: NEW-U08-R04
% key: KEY-NEW-U08-R04
% mdparagraph:P207:begin
\textbf{R04. Trade and destination.}
% mdparagraph:P207:end

% mdparagraph:P208:begin
The merchant boards his ship with his cargo and rows across stretches of sea. He sells his goods (\textbf{cīepe mīn ðing}) and buys valuable things (\textbf{bycge ðing dēorwierðu}) that are not produced in this land. He brings that merchandise \textbf{here to you}, the people he addresses (\textbf{tōgelǣde ēow hider}).
% mdparagraph:P208:end

% mdparagraph:P209:begin
\textbf{Target:} buying explains how he obtains the valuable goods; delivery explains their destination. The text does not say that he manufactures them himself. Risk during the crossing, shipwreck, loss of goods and barely escaping alive are the dangers of this activity, not additional ways of acquiring the cargo. \textbf{Ofer sǣ} means “across the sea” here. Its noun has identical accusative and dative singular forms: an accusative extent analysis is contextually motivated, but the written \textbf{sǣ} alone does not prove a unique case, and this task does not demand one.
% mdparagraph:P209:end

\end{minipage}\par\addvspace{4pt}

\par\begin{minipage}{\linewidth}\raggedright
% item: NEW-U08-R05
% key: KEY-NEW-U08-R05
% mdparagraph:P210:begin
\textbf{R05. Dangers and survival.}
% mdparagraph:P210:end

% mdparagraph:P211:begin
Any two of these expressions, correctly explained, answer the first part:
% mdparagraph:P211:end

\begin{itemize}
% mdparagraph:P212:begin
\item \textbf{Mid miclum plihte} — at great risk; the crossing puts the merchant in danger.
% mdparagraph:P212:end
% mdparagraph:P213:begin
\item \textbf{Forlidennesse ic ðolie} — I suffer shipwreck; \textbf{hwīlum} makes this an occasional experience.
% mdparagraph:P213:end
% mdparagraph:P214:begin
\item \textbf{Mid lyre ealra ðinga mīnra} — with the loss of all my things; the danger includes the loss of his merchandise or possessions.
% mdparagraph:P214:end
\end{itemize}

% mdparagraph:P215:begin
\textbf{Unēaðe cwic ætberstende} means “barely escaping alive.” \textbf{Unēaðe} gives the difficulty; \textbf{cwic}, nominative masculine singular, describes the merchant as alive; supplied present participle \textbf{ætberstende} expresses his accompanying escape. This adds survival by a narrow margin to the account of loss. “Escaping with difficulty, alive” preserves the meaning. Translating \textbf{unēaðe} as “easily,” taking \textbf{cwic} as a description of the cargo, or making the escape a passive action would lose the supplied construction.
% mdparagraph:P215:end

\end{minipage}\par\addvspace{4pt}

\reserve{7\baselineskip}
% mdparagraph:P216:begin
\section*{X. Questions, reference and clause relations}
% mdparagraph:P216:end

\par\begin{minipage}{\linewidth}\raggedright
% item: NEW-U08-X01
% key: KEY-NEW-U08-X01
% mdparagraph:P217:begin
\textbf{X01. Which men have the ships?}
% mdparagraph:P217:end

% mdparagraph:P218:begin
\textbf{Hwilce męnn habbað ðā scipu, ðe ðū gesāwe?} — Which men have the ships that you saw?
% mdparagraph:P218:end

% mdparagraph:P219:begin
\textbf{Target:} \textbf{hwilce męnn} is nominative masculine plural, the interrogative subject of \textbf{habbað}. Antecedent \textbf{ðā scipu} is accusative neuter plural, its object. Within the relative, invariant \textbf{ðe} is the object of \textbf{gesāwe}, while singular \textbf{ðū} is the subject and determines the second-singular preterite. Neither the plural men nor the plural ships makes that relative subject plural. The question asks which men, not which ships.
% mdparagraph:P219:end

\end{minipage}\par\addvspace{4pt}

\par\begin{minipage}{\linewidth}\raggedright
% item: NEW-U08-X02
% key: KEY-NEW-U08-X02
% mdparagraph:P220:begin
\textbf{X02. Another fish to the same man.}
% mdparagraph:P220:end

% mdparagraph:P221:begin
\textbf{Ic sęlle ōðerne fisc ðǣm ilcan męnn, se ðe mē healp.} — I give another fish to the same man who helped me.
% mdparagraph:P221:end

% mdparagraph:P222:begin
\textbf{Target:} \textbf{ōðerne fisc} is accusative masculine singular, with strong \textbf{ōðer} in its “another” sense. \textbf{Ðǣm ilcan męnn} is the dative recipient, with weak \textbf{ilcan}. Relative \textbf{se} is nominative masculine singular, subject of \textbf{healp}, followed by \textbf{ðe} as expressly requested; \textbf{mē} is the dative person helped. An invariant relative alone could be grammatical but would miss the required inflected member. Making \textbf{se} dative by copying the recipient's case misses the selected clause analysis.
% mdparagraph:P222:end

\end{minipage}\par\addvspace{4pt}

\par\begin{minipage}{\linewidth}\raggedright
% item: NEW-U08-X03
% key: KEY-NEW-U08-X03
% mdparagraph:P223:begin
\textbf{X03. Any little birds?}
% mdparagraph:P223:end

% mdparagraph:P224:begin
\textbf{Siehst ðū ǣnige lȳtle fuglas?} — Do you see any little birds?
% mdparagraph:P224:end

% mdparagraph:P225:begin
\textbf{Target:} second-singular present \textbf{siehst} and accusative masculine plural \textbf{ǣnige lȳtle fuglas}. \textbf{Lȳtle} has strong agreement in this indefinite phrase. \textbf{Ǣnigne fisc} was singular in Q06; the plural birds here require \textbf{ǣnige}, not that masculine singular \textbf{-ne} ending. Weak \textbf{lȳtlan} belongs to a different definite environment. An added English-style question auxiliary is unnecessary.
% mdparagraph:P225:end

\end{minipage}\par\addvspace{4pt}

\par\begin{minipage}{\linewidth}\raggedright
% item: NEW-U08-X04
% key: KEY-NEW-U08-X04
% mdparagraph:P226:begin
\textbf{X04. Each companion himself.}
% mdparagraph:P226:end

% mdparagraph:P227:begin
\textbf{Gehwilc gefēra self geseah ðæt scip.} — Each companion himself saw that ship.
% mdparagraph:P227:end

% mdparagraph:P228:begin
\textbf{Geseah gehwilc gefēra self ðæt scip?} — Did each companion himself see that ship?
% mdparagraph:P228:end

% mdparagraph:P229:begin
\textbf{Target:} singular masculine \textbf{gehwilc gefēra}, strong nominative masculine singular \textbf{self} emphasizing that subject, and third-singular preterite \textbf{geseah}. Making the finite verb initial forms the model yes/\allowbreak{}no question without changing its tense or person. \textbf{Selfne} would be an accusative form and would not agree with this subject; the object remains \textbf{ðæt scip}, not the companion himself.
% mdparagraph:P229:end

\end{minipage}\par\addvspace{4pt}

\par\begin{minipage}{\linewidth}\raggedright
% item: NEW-U08-X05
% key: KEY-NEW-U08-X05
% mdparagraph:P230:begin
\textbf{X05. A relative and a quoted jussive.}
% mdparagraph:P230:end

% mdparagraph:P231:begin
\textbf{Se mann, ðe ðæt scip hæfð, cwið: “Ætēowie drȳgnes.”} — The man who has the ship says, “Let dry land appear.”
% mdparagraph:P231:end

% mdparagraph:P232:begin
\textbf{Target:} relative \textbf{ðe} is subject of third-singular \textbf{hæfð}; \textbf{ðæt scip} is the possessed object. Main-clause \textbf{cwið} says what the man does, while the direct quotation has third-singular jussive subjunctive \textbf{ætēowie} with subject \textbf{drȳgnes}. Third-singular \textbf{sægð}, from previously taught \textbf{sęcgan}, is also valid for “says.” Retain the requested direct quotation rather than changing it to an unassigned indirect content construction. Indicative \textbf{ætēowað} would change “Let … appear” to “appears.”
% mdparagraph:P232:end

\end{minipage}\par\addvspace{4pt}

\par\begin{minipage}{\linewidth}\raggedright
% item: NEW-U08-X06
% key: KEY-NEW-U08-X06
% mdparagraph:P233:begin
\textbf{X06. Two uses of ðæt.}
% mdparagraph:P233:end

% mdparagraph:P234:begin
In \textbf{ðæt scip}, \textbf{ðæt} is an agreeing demonstrative/\allowbreak{}article with neuter singular \textbf{scip}. In the account's phrases governed by “have” or “see,” that noun phrase is accusative. It identifies the ship as “the/\allowbreak{}that ship.”
% mdparagraph:P234:end

% mdparagraph:P235:begin
In \textbf{ðæt ðe ðū gesāwe}, \textbf{ðæt} is the inflected relative member, accusative neuter singular because the ship is the thing \textbf{ðū} saw; invariant \textbf{ðe} follows it. Its antecedent is the little ship in the main clause. The whole relative means “which/\allowbreak{}that you saw.”
% mdparagraph:P235:end

% mdparagraph:P236:begin
\textbf{Target:} identical spelling does not give the two occurrences an identical grammatical function, and neither is the content conjunction in these examples. The first determines a noun; the second refers back to an antecedent and fills the object role within a relative. The separately supplied content \textbf{ðæt} in the merchant's \textbf{Ic sęcge ðæt …} is another use, not a reason to label every \textbf{ðæt} a conjunction.
% mdparagraph:P236:end

\end{minipage}\par\addvspace{4pt}

\clearpage
% mdparagraph:P237:begin
\section*{Source and editorial notes}
% mdparagraph:P237:end

\begingroup\fontsize{10}{12.8}\selectfont
% mdparagraph:P238:begin
[1] The models and explanations are new teaching material. The selected pronouns and relative strategies follow Cook, \emph{A First Book in Old English}, second edition (1895), §§86–89, with the strong and weak endings in §§57–60; supporting checks include Wright §§463–464 and 468–471 and Fulk §§35 and 77. The ordinary internal-case target is qualified by the attested recognition exception of case attraction. The lexical and verbal systems retain the lesson's selected forms and limits.
% mdparagraph:P238:end

% mdparagraph:P239:begin
[2] Cook, \emph{Exercises in Old English} (1895), I.8, printed p. 4 /\allowbreak{} PDF p. 12, all seventeen items. The English prompts are exact source material; the Old English models are new answers. Cook, \emph{First Book}, §168.1 supplies the cognate-accusative construction. The appearance witness is on printed p. 124 /\allowbreak{} PDF p. 140, with the dictionary entry \textbf{ætēowian} on p. 275 /\allowbreak{} PDF p. 291. The supplied placement and taught control frames have witnesses on printed p. 125 /\allowbreak{} PDF p. 141, line 15, and p. 127 /\allowbreak{} PDF p. 143, line 9 respectively.
% mdparagraph:P239:end

% mdparagraph:P240:begin
[3] The merchant excerpt is exact source material from Cook, \emph{First Book}, printed pp. 129–130 /\allowbreak{} PDF pp. 145–146: the three unnumbered dialogue lines on p. 129 are counted editorially; p. 130 uses printed lines 1–7. The translation covers the four selected turns through \textbf{unēaðe cwic ætberstende}. The local helps, rather than an assumed complete lexical or grammatical system, support the additional forms and constructions. Cook's original notes remain historical source paratext in the learner chapter; their antiquarian questions are not additional assigned exercises here. PDF page numbers in these notes are one-based.
% mdparagraph:P240:end

\endgroup
\end{document}
